\begin{abstract}
We compare three decoders for the square GKP code under temporally correlated
(AR(1)) displacement noise with Gaussian measurement noise, in a validated
simplified simulator of the position quadrature:
\begin{itemize}
  \item a memoryless nearest-lattice (NL) decoder;
  \item a sequential Bayesian grid filter;
  \item an LSTM trained on simulated syndrome histories.
\end{itemize}
Restricting corrections to the syndrome plus a lattice shift makes the
syndromes independent of the decoder. Decoding then reduces to estimating a
running parity, which allows paired comparisons on common noise and
supervised training without distribution shift.

At the nominal point ($\sigma_q=0.2L$, $\sigma_m=0.05L$, $\rho=0.9$), matched
Bayes reduces the logical error rate after 50 rounds from $\EoneNLLERfifty$
(NL) to $\EoneBayesLERfifty$. The LSTM trained at that point reaches
$\EoneLSTMnomLERfifty$. Being able to revise past decisions lowers the time-averaged
error by a factor of about 3 (LSTM) to 7 (Bayes).

The matched filter is the best decoder at every test point, and a moderately
mis-specified filter usually stays ahead of the LSTMs. The exception is weak
correlation ($\rho=0.3$--$0.5$), where the broadly trained LSTM beats the
model-averaged filter. Matched LSTMs come within 2--17\% of the filter where
errors are common, and fall behind by 3--4 times where they are rare. Wrong assumptions hurt both: assuming strong correlation
when there is none is worse than having no memory. Under regime-switching
noise, a broadly trained LSTM beats every fixed-model filter, and only a
regime-aware filter does better. We also report generalisation across noise
strength, correlation and time, sensitivity to model mismatch, and compute
cost.
\end{abstract}
